\documentclass[11pt]{article}
\usepackage[margin=1in]{geometry}
\usepackage{amsmath,amssymb}
\usepackage{microtype}
\usepackage[hidelinks]{hyperref}

\newcommand{\R}{\mathbb R}
\newcommand{\Sclass}{\mathcal S}

\title{Literature resolution of the Wigderson--Wigderson\\Fourier norm-ratio conjecture}
\author{Status note for run \texttt{fa\_banach\_001}}
\date{13 August 2026}

\begin{document}
\maketitle

\begin{center}
\fbox{\parbox{0.88\textwidth}{
\textbf{Status: literature already answered.}  Huang--Liu--Wu explicitly
identify Conjecture~4.13 of Wigderson--Wigderson and give its complete
truth-by-exponent resolution.
}}
\end{center}

\section*{Original conjecture}

The source is Avi Wigderson and Yuval Wigderson, \emph{The uncertainty
principle: variations on a theme}, arXiv:2006.11206 \cite{WW}.  In the current
source PDF, Conjecture~4.13 appears on PDF page~33.  For $q\in(1,\infty]$ and
$0\neq f\in\Sclass(\R)$, it defines
\[
 F_q(f)=\frac{\|f\|_q\|\widehat f\|_q}
              {\|f\|_2\|\widehat f\|_2}
       =\frac{\|f\|_q\|\widehat f\|_q}{\|f\|_2^2}
\]
and conjectures that, for every $q\neq2$, the image of $F_q$ is all of
$\R_{>0}$.  The source itself proves the endpoint $q=\infty$ in
Theorem~4.12, but leaves all finite $q\neq2$ as the conjecture.

\section*{Exact later answer}

Linzhe Huang, Zhengwei Liu, and Jinsong Wu, \emph{Quantum smooth uncertainty
principles for von Neumann bi-algebras}, arXiv:2107.09057, explicitly say that
they give a complete answer to this conjecture \cite{HLW}.  In the arXiv PDF,
Section~4 begins on page~15, restates the source statement verbatim as
``Conjecture~1 (Conjecture~4.13 in [22]),'' and Theorem~4.3 on PDF page~16
proves:
\begin{enumerate}
\item If $1<q<2$ and $p=q/(q-1)$, then every nonzero Schwartz function
satisfies
\[
 F_q(f)\geq\left(\frac{p^{1/p}}{q^{1/q}}\right)^{1/2}>0.
\]
Thus the proposed image $\R_{>0}$ is impossible in this regime.
\item If $q>2$, then the image of $F_q$ is exactly $\R_{>0}$.
\end{enumerate}
Together with the source's $q=\infty$ result, this completely resolves the
conjecture: it is true above $2$ and false below $2$.  The authors knew they
were answering the exact conjecture, rather than merely proving a theorem
that implies it after an agent-supplied identification.

The answer paper was published as \emph{Quantum Topology} 15 (2024), no.~3/4,
473--501, DOI \href{https://doi.org/10.4171/QT/207}{10.4171/QT/207}.

\section*{Scope and remaining question}

No part of Conjecture~4.13 remains open.  Huang--Liu--Wu end the section with
a new problem: for $1<q<2$, determine the optimal constant
\[
 C_q=\inf\{F_q(f):0\neq f\in\Sclass(\R)\}.
\]
That sharp-constant problem is distinct from the resolved image conjecture.

\begin{thebibliography}{9}
\bibitem{WW}
A.~Wigderson and Y.~Wigderson,
\emph{The uncertainty principle: variations on a theme},
Bull. Amer. Math. Soc. 58 (2021), 225--261;
\href{https://arxiv.org/abs/2006.11206}{arXiv:2006.11206}.

\bibitem{HLW}
L.~Huang, Z.~Liu, and J.~Wu,
\emph{Quantum smooth uncertainty principles for von Neumann bi-algebras},
Quantum Topol. 15 (2024), no.~3/4, 473--501;
\href{https://arxiv.org/abs/2107.09057}{arXiv:2107.09057}.
\end{thebibliography}

\end{document}

